\documentclass[11pt]{article}

\usepackage[margin=1in]{geometry}
\usepackage{amsmath,amssymb,amsthm}
\usepackage{graphicx}
\usepackage[hidelinks]{hyperref}

\newtheorem{theorem}{Theorem}
\newtheorem{lemma}{Lemma}
\newtheorem{remark}{Remark}

\newcommand{\N}{\mathbb N}
\newcommand{\rad}{\operatorname{rad}}
\newcommand{\linf}{\ell_\infty}
\newcommand{\norm}[1]{\left\lVert #1 \right\rVert}

\title{A Closed Bounded Radius Counterexample for \texorpdfstring{$c_0\subset \ell_\infty$}{c0 in ell-infinity}}
\author{Automated Math Discovery Smoke Test}
\date{June 8, 2026}

\begin{document}
\maketitle

\section*{Source and Target}

This packet concerns the closed-bounded version of Problem 1.1 discussed in:

\begin{quote}\raggedright
Syamantak Das and Tanmoy Paul,
\emph{A study on various generalizations of Generalized centers $(GC)$ in
Banach spaces}, arXiv:2311.15818.
\end{quote}

The paper answers the finite-subset version of Problem 1.1. In the observations
section, the authors say that the corresponding question for closed bounded
subsets is not known except in some trivial cases. The source location is page
6 of the paper.

\begin{center}
\includegraphics[width=\textwidth]{figures/open_problem_crop.png}
\end{center}

\section*{Counterexample}

Let $X=\ell_\infty$ and let $Y=c_0$ be its canonical closed subspace. For a
bounded nonempty set $A$ in a normed space $E$, write
\[
  \rad_E(A)=\inf_{x\in E}\sup_{a\in A}\norm{x-a}.
\]

\begin{theorem}
The pair $c_0\subset\ell_\infty$ has
\[
  \rad_{c_0}(F)=\rad_{\ell_\infty}(F)
\]
for every finite subset $F\subset c_0$. Nevertheless there is a closed bounded
subset $A\subset c_0$ such that
\[
  \rad_{c_0}(A)\ne\rad_{\ell_\infty}(A).
\]
In fact one may take $A=\{e_n:n\in\N\}$, for which
\[
  \rad_{\ell_\infty}(A)=\frac12,
  \qquad
  \rad_{c_0}(A)=1.
\]
\end{theorem}

\begin{proof}
We first prove the finite-subset equality. The inequality
$\rad_{\ell_\infty}(F)\le\rad_{c_0}(F)$ is immediate because
$c_0\subset\ell_\infty$. Conversely, let $F\subset c_0$ be finite and let
$r>\rad_{\ell_\infty}(F)$. Choose $z\in\ell_\infty$ such that
\[
  \sup_{a\in F}\norm{z-a}_\infty<r.
\]
Since $F$ is finite and each $a\in F$ tends to zero coordinatewise, there is
$N$ such that
\[
  \sup_{a\in F}\sup_{n>N}|a_n|<r.
\]
Define $y\in c_0$ by
\[
  y_n=
  \begin{cases}
    z_n, & n\le N,\\
    0, & n>N.
  \end{cases}
\]
For $n\le N$ the coordinate errors $|y_n-a_n|$ are the same as for $z$, and
for $n>N$ they are $|a_n|<r$. Hence
\[
  \sup_{a\in F}\norm{y-a}_\infty\le r.
\]
Letting $r\downarrow\rad_{\ell_\infty}(F)$ gives
$\rad_{c_0}(F)\le\rad_{\ell_\infty}(F)$, and the finite-subset equality
follows.

Now set $A=\{e_n:n\in\N\}$. This set is bounded in $c_0$ and is closed, since
distinct elements of $A$ are pairwise distance $1$ apart.

In $\ell_\infty$, the constant sequence $u=(1/2,1/2,\ldots)$ satisfies
\[
  \norm{u-e_n}_\infty=\frac12
  \qquad(n\in\N),
\]
so $\rad_{\ell_\infty}(A)\le1/2$. For the reverse inequality, let
$z\in\ell_\infty$ and suppose
\[
  \sup_n\norm{z-e_n}_\infty\le R.
\]
Fix $k\in\N$. From the distance to $e_k$ we get $|1-z_k|\le R$. From the
distance to any $e_m$ with $m\ne k$ we get $|z_k|\le R$. Thus
\[
  R\ge\max\{|z_k|,|1-z_k|\}\ge\frac12.
\]
Since $k$ was arbitrary, $\rad_{\ell_\infty}(A)\ge1/2$, and therefore
$\rad_{\ell_\infty}(A)=1/2$.

In $c_0$, let $y\in c_0$. Since $y_n\to0$,
\[
  \sup_n\norm{y-e_n}_\infty
  \ge
  \sup_n |1-y_n|
  \ge
  \limsup_n |1-y_n|
  =
  1.
\]
Thus $\rad_{c_0}(A)\ge1$. The center $0\in c_0$ gives
$\sup_n\norm{0-e_n}_\infty=1$, so $\rad_{c_0}(A)=1$.
\end{proof}

\section*{Interpretation}

The example shows that the finite criterion in Problem 1.1, and equivalently
the finite intersection behavior captured by $n.X.I.P.$ for all $n$, does not
control Chebyshev radii of arbitrary closed bounded subsets. The obstruction is
not finite-dimensional: every finite part of $\{e_n\}$ has an
$\ell_\infty$-center that can be truncated into $c_0$, but the limiting center
$(1/2,1/2,\ldots)$ is not in $c_0$.

\section*{Sanity Check Script}

The file \texttt{code/check\_radii.py} prints finite-prefix radius checks and
illustrates how $c_0$-like centers approach radius $1$ on longer prefixes. The
mathematical proof above is exact and does not depend on the script.

\end{document}

